% !TeX root = ../../Book1.tex
% 纲要标题：## 附录 A　集合论与拓扑预备知识
% 纲要位置：计划纲要.md，第 774 行。
% BEGIN APPENDIX OUTLINE SYNC
% * Zorn 引理与选择公理的使用清单
% * 商集、滤过系统与逆极限的集合构造
% * 积拓扑、紧性及投射有限空间的最小背景
% * 不把完整点集拓扑列为所有基础代数章节的先修
% END APPENDIX OUTLINE SYNC

\chapter{集合论与拓扑预备知识}
\label{b1:app:A}

\begin{chapterabstract}
代数中的商对象、极大对象和无限相容族，都先是集合上的构造。本附录说明这些构造用到了哪些集合论原则，整理无限 Galois 理论所需的拓扑，并补充有限支集计数、集合的正向极限和有限逆系统的基本结论。选择公理及其等价形式在这里明确列为基础约定；主干中的极大理想、代数闭包和 Krull 紧性证明仍保留在正文各自的位置。
\end{chapterabstract}

\begin{learninggoals}
\begin{enumerate}
\item 区分极大元、最大元与上界，核对 Zorn 引理中的链条件。
\item 用商集或相容族分别构造正向极限和逆极限，并验证相应的映射性质。
\item 把有限交性质、闭子空间和有限坐标条件用于投射有限空间。
\end{enumerate}
\end{learninggoals}

阅读本附录只需集合、映射和有限线性组合的基本语言。涉及理想、超越基和 Galois 群的段落是对正文的使用清单；尚未读到相应章节时可以跳过。所有指标均为集合，本附录不处理以全部集合或全部群为指标的真类构造。

% 纲要内容（仅作写作计划，尚非教材正文）：
%
% * Zorn 引理与选择公理的使用清单
% * 商集、滤过系统与逆极限的集合构造
% * 积拓扑、紧性及投射有限空间的最小背景
% * 不把完整点集拓扑列为所有基础代数章节的先修
%

% !TeX root = ../../Book1.tex
% 本节标题据《计划纲要.md》附录 A 第 1 项归纳；原要点保留于 AppendixA.tex。

\section{Zorn 引理与选择公理}
\label{b1:sec:A01}

% TODO: 按附录入口保留的纲要要点撰写本节。

待撰写。

% !TeX root = ../../Book1.tex
% 纲要标题：### A.2 商集、滤过系统与逆极限
% 纲要位置：计划纲要.md，第 789 行。
% BEGIN APPENDIX OUTLINE SYNC
% #### A.2.1 商集与映射的下降
%
% #### A.2.2 滤过系统与正向极限
%
% #### A.2.3 逆系统与相容族
% END APPENDIX OUTLINE SYNC

\section{商集、滤过系统与逆极限}
\label{b1:sec:A02}

正向系统把对象送到更大的层，逆系统则把较精细的数据限制到较粗的层。这两种构造不能只凭箭头名称记忆：前者把不同层的元素按最终相等识别，后者要求一个元素同时携带全部层的相容坐标。

\subsection{商集与映射的下降}
\label{b1:subsec:A02:01}

集合 \(X\) 上的等价关系 \(\sim\) 满足自反、对称和传递三项条件，其等价类为 \([x]=\{y\in X\mid y\sim x\}\)，商集 \(X/{\sim}\) 由这些等价类组成。\cref{b1:prop:partition-equivalence}已证明它们确实构成 \(X\) 的划分。商映射 \(q:x\mapsto[x]\) 的构造不需要先为各类选取代表元。

\begin{proposition}\label{b1:prop:set-quotient-universal}
给定集合映射 \(f:X\rightarrow Y\)，存在唯一映射 \(\bar f:X/{\sim}\rightarrow Y\) 满足 \(\bar f q=f\)，当且仅当 \(x\sim x'\) 蕴含 \(f(x)=f(x')\)。若条件成立，则 \(\bar f\) 满射当且仅当 \(f\) 满射；\(\bar f\) 单射当且仅当 \(f(x)=f(x')\) 还蕴含 \(x\sim x'\)。
\end{proposition}
\begin{proof}
若有分解，等价元素经 \(q\) 后相同，故经 \(f\) 后相同。反之，规定 \(\bar f([x])=f(x)\)；假设恰好保证它不依赖代表元。每个类都含有代表元，因而分解等式唯一确定 \(\bar f\)。两映射的像相同，给出满射判据；\(\bar f([x])=\bar f([x'])\) 等价于 \(f(x)=f(x')\)，给出单射判据。
\end{proof}
例如在整数对 \((a,b)\)、\(b\ne0\) 中按 \(ad=bc\) 识别 \((a,b)\) 与 \((c,d)\)，有理数映射 \((a,b)\mapsto a/b\) 在每类上恒定。它下降后的单射性正是“分式相等只有这一种来源”。在一般整环中，\cref{b1:prop:fraction-field-basic}核查此等价关系和运算；集合层面的映射判据并不免除这些代数核查。

\subsection{滤过系统与正向极限}
\label{b1:subsec:A02:02}

\begin{definition}\label{b1:def:directed-system}
非空偏序集 \(I\) 若任意 \(i,j\in I\) 都有共同上界，称为\term[youxiang-ji]{有向集}[directed set]。在 \(I\) 上给出集合 \(X_i\) 及映射 \(u_{ij}:X_i\rightarrow X_j\)（\(i\leq j\)），满足
\[
u_{ii}=\operatorname{id}_{X_i},\qquad
u_{jk}u_{ij}=u_{ik}\quad(i\leq j\leq k),
\]
便得到一个由有向集指标的\term[lvguo-xitong]{滤过系统}[filtered system]。这里的有向性指任意有限多个指标可放入一个共同后层，不要求任意一条无限链都有上界。
\end{definition}
任意集合的有限子集按包含构成有向集。域扩张的有限生成中间域也有向，因为两组有限生成元的并仍有限；若扩张代数，这些中间域都是有限扩张。\cref{b1:prop:finite-galois-layers}所用的有限 Galois 子扩张是带有额外正规、可分条件的同类例子。

\begin{definition}\label{b1:def:filtered-colimit-set}
对上述系统，在不交并 \(\coprod_{i\in I}X_i\) 上规定
\[
(i,x)\sim(j,y)
\quad\Longleftrightarrow\quad
\text{存在 }k\geq i,j\text{ 使 }u_{ik}(x)=u_{jk}(y).
\]
相应商集称为该系统的\term[zhengxiang-jixian]{正向极限}[direct limit]，也称\term[lvguo-yujixian]{滤过余极限}[filtered colimit]，记为 \(\varinjlim_iX_i\)。结构映射为 \(\iota_i(x)=[(i,x)]\)。
\end{definition}
\symidx{\(\varinjlim_iX_i\)：集合滤过系统的正向极限}

\begin{proposition}\label{b1:prop:filtered-colimit-universal}
上述关系是等价关系，且 \(\iota_j u_{ij}=\iota_i\)。若映射 \(f_i:X_i\rightarrow Y\) 满足 \(f_j u_{ij}=f_i\)，则存在唯一 \(f:\varinjlim_iX_i\rightarrow Y\)，使 \(f\iota_i=f_i\)。若全部 \(u_{ij}\) 单射，则每个 \(\iota_i\) 也单射；当各层原本就是同一集合中的子集且过渡映射为包含时，极限与 \(\bigcup_iX_i\) 自然双射。
\end{proposition}
\begin{proof}
自反性由 \(u_{ii}=\operatorname{id}\) 得到，对称性由等号得到。若 \((i,x)\sim(j,y)\) 的等式在层 \(k\) 成立，\((j,y)\sim(\ell,z)\) 的等式在层 \(k'\) 成立，取 \(m\geq k,k'\)。将两式映到 \(m\) 层，复合条件给出 \(u_{im}(x)=u_{jm}(y)=u_{\ell m}(z)\)，从而得到传递性。结构映射相容也直接来自定义。

规定 \(f([(i,x)])=f_i(x)\)。若两个代表元等价，在某个共同后层相等；映到 \(Y\) 后，由 \(f_k u_{ik}=f_i\) 得它们取值相同，所以良定义。所有元素都由某个 \(\iota_i(x)\) 表示，故此公式也给出唯一性。若 \(\iota_i(x)=\iota_i(y)\)，则某个 \(k\geq i\) 有 \(u_{ik}(x)=u_{ik}(y)\)；\(u_{ik}\) 单射给出 \(x=y\)。最后，在包含映射的情形，不同层元素等价恰好表示它们在母集合中相同，故 \([(i,x)]\mapsto x\) 给出与并的双射。
\end{proof}
递增域的并因而是这项集合构造的特殊情形。对相容的域嵌入，有限多个加、减、乘、除运算均可放入共同后层计算；后层映射保持运算使结果与所选层无关。若过渡映射不单射，某层两个不同元素可能在以后被识别，不能把正向极限直接画成互相包含的集合。

\subsection{逆系统与相容族}
\label{b1:subsec:A02:03}

\begin{definition}\label{b1:def:inverse-limit-set}
设 \(I\) 为有向集。一个集合的\term[nixi-tong]{逆系统}[inverse system]由集合 \(X_i\) 和映射 \(r_{ji}:X_j\rightarrow X_i\)（\(i\leq j\)）组成，满足 \(r_{ii}=\operatorname{id}\) 与 \(r_{ki}=r_{ji}r_{kj}\)。其逆极限为
\[
X\coloneqq\varprojlim_iX_i
=\Bigl\{\,(x_i)_i\in\prod_iX_i\mid
 r_{ji}(x_j)=x_i\text{ 对所有 }i\leq j\,\Bigr\}.
\]
记 \(\pi_i:X\rightarrow X_i\) 为坐标投影。这里的元素是全部坐标构成的相容族，不是某一个足够大的层中的元素。
\end{definition}

\begin{proposition}\label{b1:prop:inverse-limit-set-universal}
若集合 \(Y\) 及映射 \(f_i:Y\rightarrow X_i\) 满足 \(r_{ji}f_j=f_i\)，则存在唯一 \(f:Y\rightarrow X\) 满足 \(\pi_i f=f_i\)。
\end{proposition}
\begin{proof}
令 \(f(y)=(f_i(y))_i\)，相容性保证此族属于 \(X\)。任何满足要求的映射在每个坐标上的值都被指定，因而等于 \(f\)。
\end{proof}
若各 \(X_i\) 为群且各 \(r_{ji}\) 为同态，逐坐标运算使 \(X\) 为群；同态保持乘积与逆元，保证运算仍留在相容族内。这正是\cref{b1:def:inverse-limit}中的构造。环的情形也可逐坐标处理，但对象类型决定应保留哪些运算。

\ProseParagraphBreak
各层非空本身不保证逆极限非空。令 \(X_n=\{m\in\mathbb N\mid m\geq n\}\)，过渡映射为包含，则相容族只能是同一个自然数出现在所有坐标上；没有自然数属于全部 \(X_n\)，所以极限为空。对有限非空集合，有向性和紧性将排除这种现象；其精确陈述及投影满射的附加条件见\cref{b1:prop:finite-inverse-limit}。

% !TeX root = ../../Book1.tex
% 纲要标题：### A.3 积拓扑、紧性与投射有限空间
% 纲要位置：计划纲要.md，第 797 行。
% BEGIN APPENDIX OUTLINE SYNC
% #### A.3.1 积拓扑与连续映射
%
% #### A.3.2 有限交性质与紧性
%
% #### A.3.3 投射有限空间与有限层检测
% END APPENDIX OUTLINE SYNC

\section{积拓扑、紧性与投射有限空间}
\label{b1:sec:A03}

无限 Galois 群中的“接近”由有限层上取值相同来表达。这里只整理与这种有限信息有关的拓扑，不以度量、序列极限或完整点集拓扑作为前提。尤其对任意指标的积空间，基本开集只限制有限多个坐标。

\subsection{积拓扑与连续映射}
\label{b1:subsec:A03:01}

\begin{definition}\label{b1:def:topology-appendix}
集合 \(X\) 上的拓扑是一族包含 \(\varnothing,X\)、对任意并及有限交封闭的子集，族中成员称为开集，开集的补集称为闭集。包含一点的开集称为该点的开邻域。映射 \(f:X\rightarrow Y\) 连续，指 \(Y\) 中每个开集的原像在 \(X\) 中开。若两个不同的点总有不相交的开邻域，称空间为 Hausdorff 空间。
\end{definition}
子集 \(A\subseteq X\) 的子空间拓扑由所有 \(A\cap U\)（\(U\) 在 \(X\) 中开）组成；它的闭集恰为 \(A\cap F\)（\(F\) 在 \(X\) 中闭）。离散拓扑把每个子集都列为开集，因而离散空间中每个子集也闭。上述约定与\subsecref{b1:subsec:2202:02}一致。

\begin{definition}\label{b1:def:product-topology-appendix}
设 \((X_i)_{i\in I}\) 是拓扑空间族，\(P=\prod_iX_i\)，\(p_i:P\rightarrow X_i\) 是投影。积拓扑的基本开集为
\[
\bigcap_{i\in F}p_i^{-1}(U_i),
\qquad F\subseteq I\text{ 有限},\quad U_i\subseteq X_i\text{ 开}.
\]
所有这种集合的任意并构成积拓扑；空的 \(F\) 给出全空间。
\end{definition}
这些集合对有限交封闭，且覆盖 \(P\)，所以它们的任意并确实满足拓扑公理。若各 \(X_i\) 离散，可进一步以规定有限多个坐标的确切取值作为基本开集。特别地，\(\{0,1\}^{\mathbb N}\) 中规定所有坐标都为零得到一个闭单点，却不是开集：只固定有限坐标时总还能改动别的坐标。

\begin{proposition}\label{b1:prop:product-continuity}
映射 \(f:Y\rightarrow\prod_iX_i\) 连续，当且仅当每个坐标映射 \(p_i f\) 连续。若每个 \(X_i\) 为 Hausdorff 空间，则其积及积的任意子空间也为 Hausdorff 空间。
\end{proposition}
\begin{proof}
每个 \(p_i\) 的开集原像是基本开集，所以连续，得到必要性。反之，基本开集在 \(f\) 下的原像是有限多个 \((p_i f)^{-1}(U_i)\) 的交，故开；任意开集是基本开集的并，其原像仍开。两个不同的积空间点在某个坐标上不同；在该坐标取不相交开邻域，其投影原像将两点分离。子空间中的分离邻域只需与子空间相交。
\end{proof}
因此，\cref{b1:prop:inverse-limit-set-universal}中的极限若采用积的子空间拓扑，且各 \(f_i\) 连续，则诱导映射也连续。这说明逆极限的拓扑无需另行猜测：它就是要求全部坐标连续所产生的拓扑。

\subsection{有限交性质与紧性}
\label{b1:subsec:A03:02}

沿用正文约定：每个开覆盖有有限子覆盖称为拟紧；拟紧且 Hausdorff 才称为紧。这一点须与把“紧”直接用作“拟紧”的文献区别。集合族 \(\mathcal F\) 具有有限交性质，指任意有限个成员的交非空，包括零个成员；空交取为全空间。

\begin{proposition}\label{b1:prop:compact-fip}
空间 \(X\) 拟紧，当且仅当每个具有有限交性质的闭集族总交非空。拟紧空间的闭子空间仍拟紧；紧空间的闭子空间仍紧。
\end{proposition}
\begin{proof}
闭集族 \((F_\alpha)\) 总交为空，等价于其开补集覆盖 \(X\)；其中某个有限交为空，等价于相应有限多个开补集已经覆盖 \(X\)。将拟紧性的有限子覆盖条件取逆否命题，正好得到所述等价性。

若 \(C\subseteq X\) 闭，则 \(C\) 中的闭集也是 \(X\) 中的闭集。其有限交非空可以在 \(X\) 中检验，由刚证的判据得到总交非空，所以 \(C\) 拟紧。Hausdorff 性对子空间成立，故第二个结论也成立。空子空间的开覆盖直接有空的有限子覆盖，同样符合结论。
\end{proof}

\begin{proposition}\label{b1:prop:compact-maps}
拟紧空间在连续映射下的像仍拟紧。Hausdorff 空间中的拟紧子空间是闭集。因此，从紧空间到 Hausdorff 空间的连续双射是同胚。
\end{proposition}
\begin{proof}
像空间的开覆盖拉回给出源空间的开覆盖，有限子覆盖再映回覆盖整个像，得到第一句。设 \(C\subseteq Y\) 拟紧，\(Y\) 为 Hausdorff 空间，\(y\notin C\)。对每个 \(c\in C\)，选不相交的开邻域 \(U_c\ni c\)、\(V_c\ni y\)。\(U_c\) 覆盖 \(C\)，取有限子覆盖 \(U_{c_1},\ldots,U_{c_r}\)。开邻域 \(V_{c_1}\cap\cdots\cap V_{c_r}\) 与 \(C\) 不交，所以 \(Y\setminus C\) 开，\(C\) 闭；\(C=\varnothing\) 时直接成立。

若 \(f:X\rightarrow Y\) 为题述连续双射，\(X\) 的闭集由\cref{b1:prop:compact-fip}拟紧，其像也拟紧，进而在 \(Y\) 中闭。这表示 \(f^{-1}\) 的每个闭集原像都闭，等价于 \(f^{-1}\) 连续，故为同胚。
\end{proof}
本卷所需的积紧性是\cref{b1:lem:finite-discrete-product-compact}：任意有限离散空间族的积是紧空间。其完整证明在\subsecref{b1:subsec:2202:04}，从有有限交性质的闭集族生成真滤子，用 Zorn 引理扩张为极大真滤子，再在每个有限坐标上取得唯一被选中的纤维。有限交判据给出总交中的点。这里汇总该结论，不把正文中的证明改成对附录的省略引用，也不在此额外假设一般的积紧性定理。

\subsection{投射有限空间与有限层检测}
\label{b1:subsec:A03:03}

\begin{definition}\label{b1:def:profinite-space}
与某个有限离散集合的有向逆极限同胚的空间，称为\term[toushe-youxian-kongjian]{投射有限空间}[profinite space]。它的拓扑是相容族在各有限层直积中的子空间拓扑。若同时具有逐坐标的群结构且过渡映射都是群同态，就得到正文的投射有限群。
\end{definition}

\begin{proposition}\label{b1:prop:finite-inverse-limit}
设 \((X_i,r_{ji})\) 是非空有向集上的有限离散集合逆系统，\(X=\varprojlim_iX_i\)。则 \(X\) 紧，且其基本开集可取为 \(\pi_i^{-1}(\{a\})\)，这些集合同时闭。因而 \(X\) 全不连通。若每个 \(X_i\) 非空，则 \(X\) 非空；若过渡映射还全都满射，则每个投影 \(\pi_i:X\rightarrow X_i\) 满射。
\end{proposition}
\begin{proof}
在 \(P=\prod_iX_i\) 中，条件 \(r_{ji}(x_j)=x_i\) 定义闭集 \(C_{ji}\)：不满足条件的每个点都由这两个坐标的取值给出完全落在补集中的基本开邻域。因此 \(X=\bigcap_{i\leq j}C_{ji}\) 闭。\cref{b1:lem:finite-discrete-product-compact}和\cref{b1:prop:compact-fip}给出 \(X\) 紧。

\(X\) 中一个基本开邻域原本规定有限个坐标 \(i_1,\ldots,i_r\)。取共同上层 \(k\)，固定点的第 \(k\) 坐标便同时固定这些坐标，所以单层纤维构成邻域基。纤维的补集是同层其他纤维之并，故它也闭。两个不同点在某坐标不同，相应开闭纤维将任何同时包含它们的子空间分离为两个非空相对开部分；因此非空连通子空间只能是单点，即全不连通。

现设各层非空。采用选择公理先在 \(P\) 中取一点作为未规定坐标的默认值。给定有限多个相容条件，取高于其中全部指标的 \(k\)，任选 \(a_k\in X_k\)。将涉及的坐标 \(i\) 赋值为 \(r_{ki}(a_k)\)，其余坐标用默认值填充。复合条件保证这有限多个相容条件同时成立。所以闭集族 \((C_{ji})\) 具有有限交性质；由积的紧性和\cref{b1:prop:compact-fip}，总交 \(X\) 非空。

最后固定 \(i\) 和 \(a\in X_i\)，在上述闭集族中加入条件 \(x_i=a\)。处理任意有限个条件时，把共同上层 \(k\) 也取在 \(i\) 之上。过渡映射满射使我们能选 \(a_k\) 满足 \(r_{ki}(a_k)=a\)，于是同一有限交论证产生 \(X\) 中第 \(i\) 坐标为 \(a\) 的点，证明投影满射。
\end{proof}
非空结论与投影满射结论的假设不同。即使每层都是 \(\{0,1\}\)，若严格后层到前层的映射恒为零，极限也只有全零族，投影并不满射。有限群的逆极限总有单位族，故其非空性不需要这个存在性论证；满射问题仍要单独处理。

\begin{proposition}\label{b1:prop:finite-value-factorization}
设 \(X=\varprojlim_iX_i\) 如上，\(D\) 为有限离散集合，\(f:X\rightarrow D\) 连续。则存在指标 \(k\) 及映射 \(\bar f:\pi_k(X)\rightarrow D\)，使 \(f=\bar f\pi_k\)。若 \(\pi_k\) 满射，\(\bar f\) 定义在整个 \(X_k\) 上。
\end{proposition}
\begin{proof}
若 \(X\) 为空，取任一指标并用空映射即可。否则，每个 \(x\in X\) 都有一个单层纤维邻域 \(U_x\)，使 \(f\) 在该邻域上恒等于 \(f(x)\)，因为 \(f^{-1}(\{f(x)\})\) 开。由\cref{b1:prop:finite-inverse-limit}的紧性，有限多个 \(U_{x_1},\ldots,U_{x_r}\) 覆盖 \(X\)。取高于它们所用指标的共同上层 \(k\)。

若 \(\pi_k(y)=\pi_k(z)\)，选一个含 \(y\) 的 \(U_{x_j}\)。相容性说明 \(y,z\) 在定义此邻域的坐标上相同，所以 \(z\) 也属于它，进而 \(f(y)=f(z)\)。由\cref{b1:prop:set-quotient-universal}，\(f\) 唯一下降到 \(\pi_k(X)\)，得到所需分解。
\end{proof}
例如令 \(X_n=\mathbb Z/2^n\mathbb Z\)、\(n\geq1\)，过渡映射为约化。相容族恰好由一列数字 \(\varepsilon_0,\varepsilon_1,\ldots\in\{0,1\}\) 描述：第 \(n\) 个坐标是
\[
\varepsilon_0+2\varepsilon_1+\cdots+2^{n-1}\varepsilon_{n-1}\pmod{2^n}.
\]
给定相容族，先从模二坐标读出 \(\varepsilon_0\)；已读出前 \(n\) 位后，相容性使模 \(2^{n+1}\) 的代表元与此前部分和之差唯一为 \(0\) 或 \(2^n\)，从而确定下一位。这证明两种数据互相唯一决定。这里的“无限和”仅指全部有限截断组成的相容族，没有使用实数中的级数收敛。连续的有限值观察只能依赖有限多个起始数字；整个空间却含有无限多个点，而且没有孤立点。

% !TeX root = ../../Book1.tex
% 纲要标题：### A.4 拓扑预备知识的阅读建议
% 纲要位置：计划纲要.md，第 805 行。
% BEGIN APPENDIX OUTLINE SYNC
% <!-- 短节：不设小节 -->
% END APPENDIX OUTLINE SYNC

\section{拓扑预备知识的阅读建议}
\label{b1:sec:A04}

群、环和有限域的基础章节不需要完整的拓扑课程。商对象的底集合可由\secref{b1:sec:A02}核对；遇到一般极大理想、超越基或域嵌入延拓时，查阅\secref{b1:sec:A01}，重点检查偏序、初始对象、链上界及极大性如何推出目标。这些章节所需的代数证明均在正文，附录只把共同的集合论语言放在一处。

\ProseParagraphBreak
阅读\chapref{b1:ch:22}时，可以先沿正文的有限 Galois 层进入逆极限，再按需要回看\secref{b1:sec:A03}。必须掌握的是开集原像、有限坐标邻域、闭集的有限交性质和紧空间到 Hausdorff 空间的连续双射判据。度量化、序列紧性和一般积紧性理论都不是该章的额外前置。本附录也没有证明“每个紧 Hausdorff 全不连通空间都能写成有限空间逆极限”的反向刻画，不能仅凭这里的一向结论使用它。

\ProseParagraphBreak
后续卷中的素谱通常只具拟紧性，未必 Hausdorff；完备化中的逆极限还涉及环、模和正合性。它们需要分别增加代数或拓扑假设，不能把有限离散集合的结论原样照搬。例如各层非空但无限时逆极限可以为空，连续双射若缺少紧源或 Hausdorff 目标也未必为同胚。查阅其他教材时，应先核对“紧”的约定、过渡映射的方向以及是否假定满射，再比较定理的名字。
% ---


\begin{chaptersummary}
极大性论证的关键是选对偏序并检查链的上界；极大对象能够达到所需目标，还要另证它不能继续扩张。商集把应当相同的表达式合并，正向极限将不同层中的元素在共同后层比较，逆极限则保留所有层上相容的坐标。

\ProseParagraphBreak
有限离散空间的逆极限之所以紧，是因为相容条件在积空间中定义闭集。紧性把有限条件的可解性提升为全部条件的同时可解性，也使一个取有限离散值的连续函数只依赖某个有限层。这些结论既解释 Krull 拓扑，也提醒我们区分各层非空、各投影满射和整体具有多少个点。
\end{chaptersummary}

\BookExercises

\begin{exercise}\label{b1:ex:A-01}
设 \(X\) 是无限集合。证明有限子集按包含关系组成的偏序集是有向的，却没有极大元。说明这为什么不与 Zorn 引理矛盾。
\end{exercise}
\begin{solution}
两个有限子集的并仍有限，故给出共同上界。每个有限子集都能加入一个尚未出现的元素，所以没有极大元。在本书采用的选择公理下，选取互异的 \(x_0,x_1,\ldots\in X\)，令 \(F_n=\{x_0,\ldots,x_n\}\)。这些集合构成链，但任何包含它们的集合都含有无限多个元素，故在此偏序集中没有上界。\cref{b1:thm:A-zorn}要求每条链有上界，有向性只要求每一对元素有共同上界。
\end{solution}

\begin{exercise}\label{b1:ex:A-02}
设 \(V\) 是域 \(K\) 上的向量空间，\(A\subseteq V\) 线性无关，\(S\subseteq V\) 且 \(A\cup S\) 生成 \(V\)。证明存在 \(B\)，使 \(A\subseteq B\subseteq A\cup S\) 且 \(B\) 是 \(V\) 的基。
\end{exercise}
\begin{solution}
考虑 \(A\cup S\) 中包含 \(A\) 的线性无关子集。它非空，链的并线性无关，因为每个线性关系只涉及有限多个向量，已落在链中的一个成员中；空链则用 \(A\) 作上界。由\cref{b1:thm:A-zorn}取极大者 \(B\)。若某个 \(s\in S\) 不在 \(B\) 的线性包中，则 \(B\cup\{s\}\) 仍线性无关：一个涉及 \(s\) 的非平凡关系若其系数非零便将 \(s\) 写入线性包，若系数为零便违反 \(B\) 的无关性。这与极大性矛盾。因此 \(B\) 生成 \(A\cup S\)，进而生成 \(V\)。
\end{solution}

\begin{exercise}\label{b1:ex:A-03}
证明从 \(\mathbb N\) 到 \(\mathbb Q\) 的有限支集函数只有可数多个；再用对角论证证明从 \(\mathbb N\) 到 \(\{0,1\}\) 的全部函数不构成可数集。
\end{exercise}
\begin{solution}
有限支集函数由有限多个指标及相应有理数确定。\(\mathbb N\times\mathbb Q\) 可数，由\cref{b1:prop:finite-support-cardinality}，它的有限序列集合可数；上述编码给出所求上界，单点支集的函数又给出无限多个不同函数。

若全部零一序列可以列成 \(a^{(0)},a^{(1)},\ldots\)，定义 \(b_n=1-a^{(n)}_n\)。序列 \(b\) 在第 \(n\) 个位置不同于第 \(n\) 个已列序列，故不在列表中，矛盾。有限支集的限制不能从前一结论中删除。
\end{solution}

\begin{exercise}\label{b1:ex:A-04}
设 \(q:X\rightarrow Y\) 为满射，规定 \(x\sim x'\) 当且仅当 \(q(x)=q(x')\)。证明 \(X/{\sim}\) 与 \(Y\) 自然双射，并说明构造这个双射是否需要为每条纤维选择一个代表元。
\end{exercise}
\begin{solution}
由等号的性质得到等价关系。\cref{b1:prop:set-quotient-universal}给出映射 \([x]\mapsto q(x)\)，它因 \(q\) 满射而满射；两个像相同恰好意味着原代表元等价，所以它也单射。逆映射把 \(y\) 送到整个纤维 \(q^{-1}(y)\)，该纤维就是一个等价类，因而不需要从纤维内选出元素。选择公理只在要求进一步构造截面 \(s:Y\rightarrow X\) 时出现。
\end{solution}

\begin{exercise}\label{b1:ex:A-05}
对每个 \(n\geq0\) 取 \(X_n=\{0,1\}\)，令 \(u_{nn}\) 为恒等映射，而 \(u_{nm}\) 在 \(n<m\) 时恒取 \(0\)。求此正向系统的正向极限，并判断结构映射是否单射。
\end{exercise}
\begin{solution}
这些映射满足复合条件。任意两个带层标记的元素 \((n,a),(m,b)\)，映到严格大于 \(n,m\) 的同一层后都为零，因而由\cref{b1:def:filtered-colimit-set}等价。极限只有一个元素，各 \(X_n\) 到极限的结构映射都不是单射。正向极限可表示为递增并，需要结构映射具有额外的单射性质。
\end{solution}

\begin{exercise}\label{b1:ex:A-06}
设 \((X_i,r_{ji})\) 是有向偏序集 \(I\) 上的逆系统。若子集 \(J\subseteq I\) 满足：每个 \(i\in I\) 都有 \(j\in J\) 使 \(i\leq j\)，证明限制坐标给出 \(\varprojlim_I X_i\cong\varprojlim_J X_j\)。这里 \(J\) 采用诱导次序。
\end{exercise}
\begin{solution}
先验证 \(J\) 有向：给 \(j_1,j_2\in J\)，先在 \(I\) 中取共同上界，再取其上的 \(J\) 中元素。给定 \(J\) 上相容族 \((x_j)\)，对 \(i\in I\) 取 \(j\in J\)、\(j\geq i\)，规定 \(x_i=r_{ji}(x_j)\)。若换用 \(j'\)，在 \(J\) 中取二者共同上界 \(k\)，则两式都等于 \(r_{ki}(x_k)\)，故定义与选择无关。若 \(i\leq i'\)，在 \(i'\) 上方取 \(j\in J\)，复合律给出 \(r_{i'i}(x_{i'})=x_i\)。此延拓与限制互逆。若各层还带拓扑且过渡映射连续，延拓的每个坐标都是一个投影后复合过渡映射，故由\cref{b1:prop:product-continuity}所得双射也是同胚。
\end{solution}

\begin{exercise}\label{b1:ex:A-07}
对每个 \(n\geq0\) 取 \(X_n=\{0,1\}\)，逆系统的严格后层到前层的映射全部恒取零。证明逆极限是单点，且每个坐标投影都不满射。与 \(X_n=\{m\in\mathbb N\mid m\geq n\}\) 采用包含映射的情形比较。
\end{exercise}
\begin{solution}
相容性迫使每个坐标都是后一层映来的零，故第一系统只有全零族，投影像为 \(\{0\}\)。第二系统的相容族必须所有坐标相同，而共同值须大于等于每个自然数，不可能存在。第一系统说明有限非空层保证逆极限非空，却不保证各投影满射；第二系统说明删去有限性后，非空层甚至不能保证非空逆极限。
\end{solution}

\begin{exercise}\label{b1:ex:A-08}
在 \(X=\{0,1\}^{\mathbb N}\) 的积拓扑中，令 \(D\) 为仅有有限多个坐标为一的序列集合。证明 \(D\) 稠密但不闭，\(X\) 没有孤立点，且 \(D\) 不是紧子空间。
\end{exercise}
\begin{solution}
每个非空基本开集只规定有限个坐标，将其余坐标填零便得到 \(D\) 中的点，所以 \(D\) 稠密。全一序列不在 \(D\) 中，故 \(D\ne X\)，稠密性遂排除闭性。任意基本邻域都留有未规定的坐标，将其中一个坐标改动便得到邻域内另一点，因此没有孤立点。由\cref{b1:prop:compact-maps}，Hausdorff 空间中的紧子集必须闭，故 \(D\) 不紧。
\end{solution}

\begin{exercise}\label{b1:ex:A-09}
设 \(X\) 是有限离散空间的有向逆极限，\(C\subseteq X\) 既开又闭。证明存在某个指标 \(i\) 及 \(A\subseteq X_i\)，使 \(C=\pi_i^{-1}(A)\)。
\end{exercise}
\begin{solution}
特征函数 \(1_C:X\rightarrow\{0,1\}\) 连续，因为两个单点的原像分别为开集 \(C\) 和 \(X\setminus C\)。\cref{b1:prop:finite-value-factorization}给出它经过 \(\pi_i(X)\) 的分解。取 \(A\) 为该像中函数值为一的元素，视为 \(X_i\) 的子集，即得结论。这里无须假定 \(\pi_i\) 满射，也不必为未出现的坐标值赋予特别意义。
\end{solution}
